\documentclass[12pt, a4paper]{article}

% ── Codificação e língua ──────────────────────────────────────────────────────
\usepackage[utf8]{inputenc}
\usepackage[T1]{fontenc}
\usepackage[brazil]{babel}

% ── Layout ────────────────────────────────────────────────────────────────────
\usepackage[top=2.5cm, bottom=2.5cm, left=3cm, right=2.5cm]{geometry}
\usepackage{setspace}
\onehalfspacing

% ── Tabelas ───────────────────────────────────────────────────────────────────
\usepackage{booktabs}
\usepackage{longtable}
\usepackage{multirow}
\usepackage{array}
\usepackage{tabularx}
\usepackage{makecell}   % \thead, \makecell

% ── Cores ─────────────────────────────────────────────────────────────────────
\usepackage[table]{xcolor}
\definecolor{azul}      {RGB}{20,  60,  120}
\definecolor{verdebg}   {RGB}{220, 240, 220}
\definecolor{verdedark} {RGB}{34,  120,  34}
\definecolor{vermbg}    {RGB}{255, 220, 220}
\definecolor{vermdark}  {RGB}{160,  30,  30}
\definecolor{amarbg}    {RGB}{255, 240, 200}
\definecolor{amardark}  {RGB}{140, 100,  10}
\definecolor{cinzahd}   {RGB}{80,   80,  80}
\definecolor{azulhd}    {RGB}{245, 249, 255}
\definecolor{azulhd2}   {RGB}{230, 238, 252}

% ── Hiperlinks ────────────────────────────────────────────────────────────────
\usepackage[colorlinks=true, linkcolor=azul, citecolor=azul, urlcolor=azul]{hyperref}

% ── Captions e floats ─────────────────────────────────────────────────────────
\usepackage{caption}
\usepackage{float}
\captionsetup{font=small, labelfont=bf}

% Permite que LaTeX estique espaçamento para evitar overfull em blocos \texttt
\setlength{\emergencystretch}{3em}

% ── Comandos de conveniência ──────────────────────────────────────────────────
\newcolumntype{C}[1]{>{\centering\arraybackslash}p{#1}}
\newcolumntype{L}[1]{>{\raggedright\arraybackslash}p{#1}}
\newcolumntype{R}[1]{>{\raggedleft\arraybackslash}p{#1}}

% Célula colorida de matriz de confusão
\newcommand{\cmTP}[1]{\cellcolor{verdebg}\color{verdedark}\textbf{TP = #1}}
\newcommand{\cmTN}[1]{\cellcolor{verdebg}\color{verdedark}\textbf{TN = #1}}
\newcommand{\cmFP}[1]{\cellcolor{vermbg}\color{vermdark}\textbf{FP = #1}}
\newcommand{\cmFN}[1]{\cellcolor{vermbg}\color{vermdark}\textbf{FN = #1}}

% Células de resultado na tabela por voo
\newcommand{\resTP} {\cellcolor{verdebg}\color{verdedark}\textbf{TP}}
\newcommand{\resFN} {\cellcolor{vermbg}\color{vermdark}\textbf{FN}}
\newcommand{\resFPa}{\cellcolor{amarbg}\color{amardark}\textbf{FP antes}}
\newcommand{\resFP} {\cellcolor{vermbg}\color{vermdark}\textbf{FP}}
\newcommand{\resTN} {\cellcolor{verdebg}\color{verdedark}\textbf{TN}}

% ─────────────────────────────────────────────────────────────────────────────
\begin{document}

\title{\textbf{Detecção de Falha de Motor em VANTs com\\Aprendizado de Máquina}}
\author{}
\date{}
\maketitle
\thispagestyle{empty}

% ══════════════════════════════════════════════════════════════════════════════
\section{Ideia central}
% ══════════════════════════════════════════════════════════════════════════════

O projeto desenvolve um sistema de detecção automática de falha de motor em voo
para uma aeronave não tripulada de asa fixa (CarbonZ). A motivação é crítica: em
voo autônomo, identificar o momento exato da perda de tração com latência mínima
é o que viabiliza uma resposta de emergência eficaz --- seja um pouso de
emergência ou um comando de autorrotação. O conjunto de dados conta com 33~voos
registrados em formato ROS, sendo 23 com falhas induzidas de motor e 10 voos
normais, com dados de múltiplos sensores coletados a ${\sim}520$~Hz.

% ══════════════════════════════════════════════════════════════════════════════
\section{Metodologia}
% ══════════════════════════════════════════════════════════════════════════════

A solução foi construída como um \textit{pipeline} Kedro completo, cobrindo desde
a ingestão dos tópicos ROS (com alinhamento temporal entre sensores via
\texttt{merge\_asof}) até o treinamento dos modelos. Foram desenvolvidas duas
trilhas paralelas de engenharia de \textit{features}: a primeira
(\texttt{all\_features}) extrai grandezas físicas como energia específica,
esforço de controle e razão de planeio, além de estatísticas em janelas
deslizantes (média, desvio-padrão e inclinação linear); a segunda
(\texttt{fft\_features}) opera no domínio da frequência, extraindo potência de
pico, entropia espectral e razão de alta frequência dos sinais do IMU e
magnetômetro. A seleção das 20 melhores \textit{features} foi feita pelo
$d$~de~Cohen, que mede a separabilidade entre as distribuições em voo normal e
em falha sem o viés de autocorrelação temporal que afetaria a importância de
\textit{Random~Forest}. O modelo escolhido foi \textit{Isolation Forest},
treinado com janelas deslizantes de 20 amostras por ponto, no regime não
supervisionado --- adequado à escassez de exemplos rotulados de falha.

% ══════════════════════════════════════════════════════════════════════════════
\section{Resultados e \textit{trade-off} entre os modelos}
% ══════════════════════════════════════════════════════════════════════════════

Os dois modelos têm perfis de erro opostos, o que revela que capturam fenômenos
físicos distintos. O modelo \texttt{all\_features} detectou 16 das 23~falhas
(69{,}6\%) com latência mediana de 3{,}15~s, reagindo às \textit{consequências}
da falha: queda de energia, divergência entre \textit{throttle} e velocidade,
acúmulo de erros de controle. Seu ponto fraco é a baixa revocação na classe de
falha (18{,}4\%) --- ele é conservador e tende a silenciar em casos ambíguos.

O modelo \texttt{fft\_features} detectou 15~falhas (65{,}2\%) com latência
mediana de 2{,}64~s, mas a maioria das detecções ocorreu \textit{antes} do
momento oficial da falha --- o que sugere que ele captura a assinatura de
vibração e caos espectral da degradação do motor como \textit{precursor}, não
como consequência. O custo é alto: 60\% dos voos normais geraram ao menos um
alarme falso no \texttt{all\_features}, contra 30\% no \texttt{fft\_features}.

Em termos de aplicação, o \texttt{all\_features} é preferível quando a
confiabilidade operacional importa mais --- cada alarme tem alta probabilidade de
ser real. Já o \texttt{fft\_features} é mais adequado em contextos de detecção
precoce tolerantes a falsos positivos, por exemplo como primeiro estágio de um
sistema em cascata onde sua detecção antecipada dispara o \texttt{all\_features}
como confirmatório. Um \textit{ensemble} com votação ponderada pode combinar a
sensibilidade espectral do FFT com a especificidade física do
\texttt{all\_features}.

% ══════════════════════════════════════════════════════════════════════════════
\section{Métricas por modelo e classe (nível de janela)}
% ══════════════════════════════════════════════════════════════════════════════

A Tabela~\ref{tab:metricas} apresenta as métricas de classificação calculadas
pela função \texttt{classifi\-cation\_re\-port} do scikit-learn sobre o conjunto
de teste (últimos 30\% de cada voo, em ordem temporal). O elevado
\textit{recall} da classe~0 (Normal) no modelo \texttt{all\_features} (0{,}999)
reflete o comportamento conservador do \textit{Isolation Forest} configurado com
\texttt{contamination\,=\,0{,}03}: apenas 53 das 93.289 janelas normais são
erroneamente classificadas como anomalia (taxa de falso positivo por janela de
0{,}057\%).

\begin{table}[H]
  \centering
  \caption{Métricas de classificação por modelo e classe (nível de janela de 20 amostras).}
  \label{tab:metricas}
  \begin{tabular}{llccc}
    \toprule
    \textbf{Modelo} & \textbf{Classe} & \textbf{Precisão} & \textbf{Recall} & \textbf{F1} \\
    \midrule
    \rowcolor{azulhd}
    \texttt{all\_features} & Normal (0) & 0{,}874 & 0{,}999 & 0{,}932 \\
    \rowcolor{azulhd}
    \texttt{all\_features} & Falha  (1) & 0{,}983 & 0{,}184 & 0{,}310 \\
    \rowcolor{azulhd2}
    \texttt{fft\_features} & Normal (0) & 0{,}856 & 0{,}955 & 0{,}903 \\
    \rowcolor{azulhd2}
    \texttt{fft\_features} & Falha  (1) & 0{,}402 & 0{,}158 & 0{,}227 \\
    \bottomrule
  \end{tabular}
\end{table}

% ══════════════════════════════════════════════════════════════════════════════
\section{Matrizes de confusão --- nível de voo}
% ══════════════════════════════════════════════════════════════════════════════

Cada voo é classificado como um único evento: se o modelo emitiu alarme durante
o intervalo de falha, o voo é TP; se nunca emitiu (ou só emitiu antes da janela
de falha), é~FN. Voos sem falha que geraram ao menos um alarme são FP; os
demais, TN.

% ─── all_features ─────────────────────────────────────────────────────────────
\subsection{\texttt{all\_features}}

\noindent
\textbf{Falhas detectadas:} 16/23 \quad
\textbf{Falhas perdidas:} 7/23 \quad
\textbf{Lat.\ média:} 6{,}72~s \quad
\textbf{Lat.\ mediana:} 3{,}15~s

\vspace{6pt}

\begin{table}[H]
  \centering
  \caption{Matriz de confusão (nível de voo) --- \texttt{all\_features}.}
  \label{tab:cm_all}
  \begin{tabular}{lcc}
    \toprule
    & \textbf{Alarme (positivo)} & \textbf{Silêncio (negativo)} \\
    \midrule
    \textbf{Real: com falha} & \cmTP{16} & \cmFN{7} \\
    \textbf{Real: sem falha} & \cmFP{6}  & \cmTN{4} \\
    \bottomrule
  \end{tabular}
\end{table}

\clearpage
\small
\begin{longtable}{L{3.4cm}C{1.5cm}C{1.7cm}C{1.5cm}C{2.0cm}C{1.7cm}}
  \caption{Detalhe por voo com falha --- \texttt{all\_features}.}
  \label{tab:voos_falha_all} \\
  \toprule
  \textbf{Voo} & \textbf{Falha (s)} & \textbf{Detec.\ (s)} & \textbf{Lat.\ (s)} & \textbf{Resultado} & \textbf{FP antes} \\
  \midrule
  \endfirsthead
  \multicolumn{6}{l}{\small\textit{(continuação da Tabela~\ref{tab:voos_falha_all})}} \\
  \toprule
  \textbf{Voo} & \textbf{Falha (s)} & \textbf{Detec.\ (s)} & \textbf{Lat.\ (s)} & \textbf{Resultado} & \textbf{FP antes} \\
  \midrule
  \endhead
  \midrule
  \multicolumn{6}{r}{\small\textit{(continua na próxima página)}} \\
  \endfoot
  \bottomrule
  \endlastfoot
  2018-07-18 \#1  & 115,3 & ---   & ---   & \resFN  & 0   \\
  2018-07-18 \#2  & 72,4  & ---   & ---   & \resFN  & 0   \\
  2018-07-18 EMR  & 115,6 & 117,6 & 1,92  & \resTP  & 0   \\
  2018-07-18 EMR2 & 113,1 & 115,8 & 2,72  & \resTP  & 0   \\
  2018-07-30 EMR  & 122,2 & 133,1 & 10,98 & \resTP  & 0   \\
  2018-07-30 \#1  & 115,8 & ---   & ---   & \resFN  & 0   \\
  2018-07-30 \#2  & 90,7  & ---   & ---   & \resFPa & 126 \\
  2018-07-30 EMR2 & 116,2 & 129,1 & 12,91 & \resTP  & 0   \\
  2018-07-30 EMR3 & 86,7  & 94,0  & 7,34  & \resTP  & 0   \\
  2018-07-30 EMR4 & 132,5 & 135,4 & 2,92  & \resTP  & 0   \\
  2018-07-30 EMR5 & 89,4  & 100,5 & 11,16 & \resTP  & 0   \\
  2018-09-11 \#1  & 102,6 & 103,3 & 0,72  & \resTP  & 0   \\
  2018-09-11 \#2  & 103,9 & ---   & ---   & \resFN  & 0   \\
  2018-09-11 \#3  & 48,9  & 59,5  & 10,62 & \resTP  & 0   \\
  2018-10-05 \#1  & 48,2  & 64,2  & 16,00 & \resTP  & 0   \\
  2018-10-05 \#2  & 99,1  & 101,7 & 2,58  & \resTP  & 0   \\
  2018-10-05 \#3  & 75,2  & 91,2  & 15,99 & \resTP  & 0   \\
  2018-10-18 \#1  & 103,2 & 106,0 & 2,76  & \resTP  & 0   \\
  2018-10-18 \#2  & 110,1 & 113,1 & 3,00  & \resTP  & 0   \\
  2018-10-18 \#3  & 99,4  & ---   & ---   & \resFN  & 0   \\
  2018-10-18 \#4  & 97,3  & ---   & ---   & \resFN  & 0   \\
  2018-10-18 \#5  & 100,3 & 103,6 & 3,29  & \resTP  & 0   \\
  2018-10-18 \#6  & 101,6 & 104,2 & 2,66  & \resTP  & 53  \\
\end{longtable}
\normalsize

\begin{table}[H]
  \centering
  \caption{Detalhe por voo sem falha --- \texttt{all\_features}.}
  \label{tab:voos_normal_all}
  \small
  \begin{tabular}{L{5cm}C{3.5cm}C{2cm}}
    \toprule
    \textbf{Voo} & \textbf{Alarmes falsos (janelas)} & \textbf{Resultado} \\
    \midrule
    2018-07-18 sem-falha    & 0   & \resTN \\
    2018-07-30 sem-falha    & 4   & \resFP \\
    2018-09-11 \#1 sem-falha & 72  & \resFP \\
    2018-09-11 \#2 sem-falha & 0   & \resTN \\
    2018-09-11 \#3 sem-falha & 162 & \resFP \\
    2018-10-05 \#1 sem-falha & 525 & \resFP \\
    2018-10-05 \#2 sem-falha & 64  & \resFP \\
    2018-10-05 \#3 sem-falha & 57  & \resFP \\
    2018-10-05 \#4 sem-falha & 0   & \resTN \\
    2018-10-18 sem-falha    & 0   & \resTN \\
    \bottomrule
  \end{tabular}
\end{table}

\clearpage
% ─── fft_features ─────────────────────────────────────────────────────────────
\subsection{\texttt{fft\_features}}

\noindent
\textbf{Falhas detectadas:} 15/23 \quad
\textbf{Falhas perdidas:} 8/23 \quad
\textbf{Lat.\ média:} 4{,}57~s \quad
\textbf{Lat.\ mediana:} 2{,}64~s

\vspace{6pt}

\begin{table}[H]
  \centering
  \caption{Matriz de confusão (nível de voo) --- \texttt{fft\_features}.}
  \label{tab:cm_fft}
  \begin{tabular}{lcc}
    \toprule
    & \textbf{Alarme (positivo)} & \textbf{Silêncio (negativo)} \\
    \midrule
    \textbf{Real: com falha} & \cmTP{15} & \cmFN{8} \\
    \textbf{Real: sem falha} & \cmFP{3}  & \cmTN{7} \\
    \bottomrule
  \end{tabular}
\end{table}

\clearpage
\small
\begin{longtable}{L{3.4cm}C{1.5cm}C{1.7cm}C{1.5cm}C{2.0cm}C{1.7cm}}
  \caption{Detalhe por voo com falha --- \texttt{fft\_features}.}
  \label{tab:voos_falha_fft} \\
  \toprule
  \textbf{Voo} & \textbf{Falha (s)} & \textbf{Detec.\ (s)} & \textbf{Lat.\ (s)} & \textbf{Resultado} & \textbf{FP antes} \\
  \midrule
  \endfirsthead
  \multicolumn{6}{l}{\small\textit{(continuação da Tabela~\ref{tab:voos_falha_fft})}} \\
  \toprule
  \textbf{Voo} & \textbf{Falha (s)} & \textbf{Detec.\ (s)} & \textbf{Lat.\ (s)} & \textbf{Resultado} & \textbf{FP antes} \\
  \midrule
  \endhead
  \midrule
  \multicolumn{6}{r}{\small\textit{(continua na próxima página)}} \\
  \endfoot
  \bottomrule
  \endlastfoot
  2018-07-18 \#1  & 115,3 & ---   & ---   & \resFN  & 69   \\
  2018-07-18 \#2  & 72,4  & ---   & ---   & \resFN  & 264  \\
  2018-07-18 EMR  & 115,6 & 117,8 & 2,17  & \resTP  & 755  \\
  2018-07-18 EMR2 & 113,1 & 115,9 & 2,85  & \resTP  & 0    \\
  2018-07-30 EMR  & 122,2 & 127,2 & 5,07  & \resTP  & 380  \\
  2018-07-30 \#1  & 115,8 & 130,4 & 14,64 & \resTP  & 543  \\
  2018-07-30 \#2  & 90,7  & ---   & ---   & \resFN  & 220  \\
  2018-07-30 EMR2 & 116,2 & ---   & ---   & \resFN  & 0    \\
  2018-07-30 EMR3 & 86,7  & 89,5  & 2,80  & \resTP  & 0    \\
  2018-07-30 EMR4 & 132,5 & 154,5 & 21,99 & \resTP  & 74   \\
  2018-07-30 EMR5 & 89,4  & 98,3  & 8,97  & \resTP  & 258  \\
  2018-09-11 \#1  & 102,6 & 104,0 & 1,41  & \resTP  & 57   \\
  2018-09-11 \#2  & 103,9 & ---   & ---   & \resFN  & 26   \\
  2018-09-11 \#3  & 48,9  & ---   & ---   & \resFN  & 35   \\
  2018-10-05 \#1  & 48,2  & 48,2  & 0,04  & \resTP  & 0    \\
  2018-10-05 \#2  & 99,1  & 99,3  & 0,19  & \resTP  & 26   \\
  2018-10-05 \#3  & 75,2  & ---   & ---   & \resFN  & 93   \\
  2018-10-18 \#1  & 103,2 & 105,9 & 2,64  & \resTP  & 320  \\
  2018-10-18 \#2  & 110,1 & ---   & ---   & \resFN  & 579  \\
  2018-10-18 \#3  & 99,4  & 99,7  & 0,27  & \resTP  & 1455 \\
  2018-10-18 \#4  & 97,3  & 100,1 & 2,80  & \resTP  & 75   \\
  2018-10-18 \#5  & 100,3 & 100,7 & 0,39  & \resTP  & 654  \\
  2018-10-18 \#6  & 101,6 & 103,8 & 2,25  & \resTP  & 693  \\
\end{longtable}
\normalsize

\begin{table}[H]
  \centering
  \caption{Detalhe por voo sem falha --- \texttt{fft\_features}.}
  \label{tab:voos_normal_fft}
  \small
  \begin{tabular}{L{5cm}C{3.5cm}C{2cm}}
    \toprule
    \textbf{Voo} & \textbf{Alarmes falsos (janelas)} & \textbf{Resultado} \\
    \midrule
    2018-07-18 sem-falha    & 0   & \resTN \\
    2018-07-30 sem-falha    & 184 & \resFP \\
    2018-09-11 \#1 sem-falha & 0   & \resTN \\
    2018-09-11 \#2 sem-falha & 278 & \resFP \\
    2018-09-11 \#3 sem-falha & 0   & \resTN \\
    2018-10-05 \#1 sem-falha & 0   & \resTN \\
    2018-10-05 \#2 sem-falha & 0   & \resTN \\
    2018-10-05 \#3 sem-falha & 332 & \resFP \\
    2018-10-05 \#4 sem-falha & 0   & \resTN \\
    2018-10-18 sem-falha    & 0   & \resTN \\
    \bottomrule
  \end{tabular}
\end{table}

\end{document}
